Write \(U=\{a\in\mathbb C:\operatorname{Re}a>0\}\).
On this half-plane, \(\operatorname{Log}\) denotes the principal logarithm,
which is real on the positive axis. All powers of \(a+x\), for \(x\geq0\),
use this logarithm. For \(\operatorname{Re}w>1\), define
\[
 \zeta_H(w,a)=\sum_{r=0}^{\infty}(a+r)^{-w},\qquad a\in U.
\]
On every compact subset of this domain, the summands are bounded by
\(C(1+r)^{-1-\varepsilon}\) for some \(C,\varepsilon>0\).
Thus the sum is jointly holomorphic in \((w,a)\).

\begin{lemma}[A summation formula with a controlled remainder]
\label{lem:osc-summation}
For real \(t\), let \(\{t\}=t-\lfloor t\rfloor\) and set
\(q(t)=\{t\}(1-\{t\})/2\). For every integer \(N\geq0\),
\begin{align}
 \zeta_H(w,a)
 &=\sum_{r=0}^{N-1}(a+r)^{-w}
   +\frac{(N+a)^{1-w}}{w-1}+\frac12(N+a)^{-w}
   +R_N(w,a),                                      \label{eq:osc-cutoff}\\
 R_N(w,a)
 &=w(w+1)\int_N^\infty q(t)(t+a)^{-w-2}\,dt.       \label{eq:osc-remainder}
\end{align}
The integral is jointly holomorphic for \(\operatorname{Re}w>-1\),
\(a\in U\). Formula~\eqref{eq:osc-cutoff} gives a meromorphic continuation
there, with only a simple pole at \(w=1\), of residue \(1\).
In particular,
\begin{equation}
 \zeta_H(0,a)=\frac12-a,\qquad a\in U.              \label{eq:osc-zeta-zero}
\end{equation}
\end{lemma}

\begin{proof}
On \([n,n+1]\), the polynomial \(q(t)=(t-n)(n+1-t)/2\)
vanishes at both endpoints. Its endpoint derivatives are \(1/2\) and
\(-1/2\), and its second derivative is \(-1\).
For any twice continuously differentiable complex function \(f\),
two integrations by parts give
\[
 \int_n^{n+1}q(t)f''(t)\,dt
 =\frac{f(n)+f(n+1)}2-\int_n^{n+1}f(t)\,dt.
\]
Sum this identity from \(n=N\) to \(n=M-1\). The endpoints give
\[
 \sum_{r=N}^{M-1}f(r)
 =\int_N^M f(t)\,dt+\frac{f(N)-f(M)}2
   +\int_N^M q(t)f''(t)\,dt.
\]
Take \(f(t)=(t+a)^{-w}\) and first assume \(\operatorname{Re}w>1\).
Then \(f''(t)=w(w+1)(t+a)^{-w-2}\).
The terms at \(M\) vanish as \(M\to\infty\), while
\(\int_N^\infty f(t)\,dt=(N+a)^{1-w}/(w-1)\).
This proves~\eqref{eq:osc-cutoff} in the initial domain.

For joint continuation, fix compact sets \(K\subset U\) and
\(W\subset\{\operatorname{Re}w>-1\}\).
Put \(\delta=\min_K\operatorname{Re}a>0\) and
\(\sigma=\min_W\operatorname{Re}w>-1\).
The bounds \(0\leq q\leq1/8\) and
\(t+\delta\leq|t+a|\leq t+\max_K|a|\) imply
\[
 |q(t)(t+a)^{-w-2}|
 \leq C_{K,W}(1+t)^{-\sigma-2}.
\]
Indeed, \(|\arg(t+a)|<\pi/2\), and \(\operatorname{Im}w\) is bounded
on \(W\). The same bound for any fixed parameter derivatives acquires
only powers of \(\log(2+t)\) and additional nonpositive powers of \(1+t\).
These bounds remain integrable. Differentiation under the integral is
therefore valid locally in both variables. They also give
\[
 |R_N(w,a)|\leq C_{K,W}|w(w+1)|(1+N)^{-\sigma-1}.
\]
All other terms in~\eqref{eq:osc-cutoff} are holomorphic except the displayed
pole at \(w=1\). Its residue is \((N+a)^0=1\).
The continued formulas for different \(N\) agree, since they agree for
\(\operatorname{Re}w>1\). At zero, the remainder vanishes, and the
remaining terms give \(N-(N+a)+1/2\).
\end{proof}

The derivative at zero involves logarithms of finite products.
The normalization is that of the Gamma integral already used in the
local real factor:
\begin{equation}
 \Gamma(a)=\int_0^\infty e^{-x}x^{a-1}\,dx,
 \qquad a\in U,                                   \label{eq:osc-gamma-integral}
\end{equation}
where \(x^{a-1}=\exp((a-1)\log x)\) uses the real logarithm.

\begin{lemma}[The gamma logarithm and its reciprocal]
\label{lem:osc-gamma}
The integral~\eqref{eq:osc-gamma-integral} is holomorphic and nonzero on
\(U\). There is a unique holomorphic logarithm \(L\) of \(\Gamma\)
which is real on the positive axis. It satisfies, locally uniformly on \(U\),
\begin{equation}
 L(a)=\lim_{N\to\infty}
 \left\{\log((N-1)!)+a\log N
              -\sum_{r=0}^{N-1}\operatorname{Log}(a+r)\right\}.
 \label{eq:osc-gamma-log-limit}
\end{equation}
The reciprocal of \(\Gamma\) extends to the entire function
\begin{equation}
 E(a)=a e^{\gamma a}\prod_{n=1}^{\infty}(1+a/n)e^{-a/n},
 \qquad
 \gamma=\lim_{N\to\infty}\left(\sum_{n=1}^N\frac1n-\log N\right).
 \label{eq:osc-gamma-product}
\end{equation}
Its zeros are simple and occur exactly at \(0,-1,-2,\ldots\).
Moreover, \(\Gamma(1)=1\) and \(\Gamma(a+1)=a\Gamma(a)\) on \(U\).
\end{lemma}

\begin{proof}
On a compact set with \(0<\delta\leq\operatorname{Re}a\leq A\),
the integral and each fixed derivative have an integrable majorant
\[
 e^{-x}(x^{\delta-1}+x^{A-1})(1+|\log x|^k).
\]
This proves holomorphy. Direct integration gives \(\Gamma(1)=1\).
Integration by parts, with vanishing endpoints, gives the recurrence.

For \(N\geq1\), repeated integration by parts on \([0,1]\) gives
\[
 \int_0^1 t^{a-1}(1-t)^N\,dt
 =\frac{N!}{a(a+1)\cdots(a+N)}.
\]
At each step, integrate \(t^{a+j-1}\) and differentiate the remaining
integer power of \(1-t\). The last integral is \(1/(a+N)\).
The substitution \(x=Nt\) yields
\[
 \frac{N!N^a}{a(a+1)\cdots(a+N)}
 =\int_0^N x^{a-1}(1-x/N)^N\,dx
 \longrightarrow\Gamma(a).
\]
For \(0\leq x<N\), one has \((1-x/N)^N\leq e^{-x}\).
Extend the integrand by zero for \(x\geq N\).
The preceding majorants prove convergence locally uniformly in \(a\),
including any fixed number of derivatives.

Let \(H_N=\sum_{n=1}^N1/n\).
The inequality \(\log(1+1/N)>1/(N+1)\) shows that \(H_N-\log N\)
decreases. Integral comparison gives \(H_N\geq\log(N+1)\), so the
sequence has a finite limit \(\gamma\).
The reciprocals of the finite expressions above are
\[
 aN^{-a}\prod_{n=1}^N(1+a/n)
 =a\exp\bigl(a(H_N-\log N)\bigr)
       \prod_{n=1}^N(1+a/n)e^{-a/n}.
\]
For \(|u|\leq1/2\), the logarithm series gives
\( |\operatorname{Log}(1+u)-u|\leq |u|^2\).
Thus the logarithms of the product tails converge uniformly on each
compact subset of \(\mathbb C\). Each tail is the exponential of a
holomorphic sum and is nonzero. The finite initial factors prove that
\(E\) is entire and identify every zero and its order.
Multiplying the two finite expressions and passing to their locally
uniform limits gives \(\Gamma E=1\) on \(U\).

On \(U\), define explicitly
\[
 L(a)=-\operatorname{Log}a-\gamma a
       -\sum_{n=1}^{\infty}
        \bigl(\operatorname{Log}(1+a/n)-a/n\bigr).
\]
The locally uniform convergence proves holomorphy, and
\(\exp L=1/E=\Gamma\). Every term is real for \(a>0\).
Two holomorphic logarithms differ by a constant in \(2\pi i\mathbb Z\),
so the stated real condition proves uniqueness.
Finally, the expression in braces in~\eqref{eq:osc-gamma-log-limit} equals
\[
 -\operatorname{Log}a+a(\log N-H_{N-1})
 -\sum_{n=1}^{N-1}
        \bigl(\operatorname{Log}(1+a/n)-a/n\bigr).
\]
Here \(H_{N-1}-\log N\to\gamma\), which proves the limit with its
specified branch.
\end{proof}

\begin{lemma}[The factorial constant]
\label{lem:osc-stirling}
As positive integers \(N\) tend to infinity,
\begin{equation}
 \log((N-1)!)
 =\left(N-\frac12\right)\log N-N
      +\frac12\log(2\pi)+o(1).                    \label{eq:osc-stirling}
\end{equation}
Also \(\Gamma(1/2)=\sqrt\pi\).
\end{lemma}

\begin{proof}
Let \(J=\int_{\mathbb R}e^{-u^2/2}\,du\), which is finite and positive.
The product integral in polar coordinates gives
\[
 J^2=\int_{\mathbb R^2}e^{-(u^2+v^2)/2}\,du\,dv
 =2\pi\int_0^\infty e^{-r^2/2}r\,dr=2\pi.
\]
Positivity gives \(J=\sqrt{2\pi}\).
Substitution \(x=u^2\) in~\eqref{eq:osc-gamma-integral}, followed by
rescaling this Gaussian integral, gives \(\Gamma(1/2)=\sqrt\pi\).

The recurrence for \(\Gamma\) gives \(n!=\int_0^\infty x^n e^{-x}\,dx\).
Set \(x=n+\sqrt n\,u\). Then
\begin{equation}
 \frac{n!}{n^n e^{-n}\sqrt n}
 =\int_{-\sqrt n}^{\infty}
   \exp\left(n\left[\log(1+u/\sqrt n)-u/\sqrt n\right]\right)du.
 \label{eq:osc-laplace}
\end{equation}
For each fixed \(u\), the exponent tends to \(-u^2/2\).
We give bounds that justify integrating this limit.
For \(-1<x\leq0\), differentiation shows
\(\log(1+x)-x\leq-x^2/2\).
For \(x\geq0\), integration of \(t/(1+t)\) gives
\[
 x-\log(1+x)=\int_0^x\frac{t}{1+t}\,dt
 \geq\frac{x^2}{2(1+x)}.
\]
Consequently the integrand in~\eqref{eq:osc-laplace} is bounded by
\(e^{-u^2/2}\) on \((-\sqrt n,0]\), and by \(e^{-u^2/4}\) on
\([0,\sqrt n]\). Extend these portions by zero to apply dominated
convergence on the two half-lines.

For the remaining tail, \(\log(1+x)/x\) decreases for \(x>0\).
Its derivative has the sign of
\(x/(1+x)-\log(1+x)<0\), as integration of \(1/(1+t)\) shows.
Thus, for \(x\geq1\),
\(\log(1+x)-x\leq-\kappa x\), where \(\kappa=1-\log2>0\).
The tail in~\eqref{eq:osc-laplace} is at most
\[
 \int_{\sqrt n}^{\infty}e^{-\kappa\sqrt n\,u}\,du
 =\frac{e^{-\kappa n}}{\kappa\sqrt n}\longrightarrow0.
\]
It follows that \(n!/(n^n e^{-n}\sqrt n)\to\sqrt{2\pi}\).
Take real logarithms, then use \((N-1)!=N!/N\), to obtain
\eqref{eq:osc-stirling}.
\end{proof}

\begin{proposition}[The two Hurwitz values]
\label{prop:osc-hurwitz-values}
The continuation in Lemma~\ref{lem:osc-summation} is holomorphic near
\(\{0\}\times U\), and
\begin{equation}
 \zeta_H(0,a)=\frac12-a,\qquad
 \partial_w\zeta_H(0,a)=L(a)-\frac12\log(2\pi).
 \label{eq:osc-hurwitz-values}
\end{equation}
Here \(L=\log\Gamma\) is exactly the holomorphic branch constructed in
Lemma~\ref{lem:osc-gamma}.
\end{proposition}

\begin{proof}
The first identity and joint holomorphy were proved in
Lemma~\ref{lem:osc-summation}. Differentiating its exact formula at zero
gives, with \(A=N+a\),
\begin{align*}
 \partial_w\zeta_H(0,a)
 ={}&-\sum_{r=0}^{N-1}\operatorname{Log}(a+r)
       +(A-\tfrac12)\operatorname{Log}A-A\\
    &+\int_N^\infty\frac{q(t)}{(t+a)^2}\,dt.
\end{align*}
The last integral has absolute value at most \(1/(8(N+\delta))\)
when \(\operatorname{Re}a\geq\delta>0\).
It therefore tends to zero locally uniformly on \(U\).
Use~\eqref{eq:osc-gamma-log-limit} to rewrite the remaining expression
as the convergent approximation to \(L(a)\), plus
\[
 -\log((N-1)!)-a\log N
 +(N+a-\tfrac12)\operatorname{Log}(N+a)-(N+a).
\]
By~\eqref{eq:osc-stirling}, this equals
\[
 -\frac12\log(2\pi)
 +(N+a-\tfrac12)\operatorname{Log}(1+a/N)-a+o(1).
\]
For \(a\) in a fixed compact set,
\(\operatorname{Log}(1+a/N)=a/N+O(N^{-2})\) uniformly.
The last two terms therefore tend to zero locally uniformly.
This proves the derivative identity for complex \(a\in U\) directly.
\end{proof}
